%==============================================================================
% preamble.tex — visual identity matching the web report
%
% Palette, type and box styles are lifted from the HTML version so the PDF and
% the artifact read as the same document. Everything here is presentation; no
% content lives in this file.
%==============================================================================

%------------------------------------------------------------------------------
% PALETTE  (same hex values as the web report's CSS custom properties)
%------------------------------------------------------------------------------
\usepackage[table,dvipsnames]{xcolor}
\definecolor{paper}{HTML}{EEF0EA}   % page ground (sage-grey)
\definecolor{raised}{HTML}{F6F7F3}  % box fill, one step lighter
\definecolor{ink}{HTML}{1B211C}     % body text, green-black
\definecolor{muted}{HTML}{5F6A62}   % captions, secondary
\definecolor{rule}{HTML}{CDD3C9}    % hairlines
\definecolor{rulesoft}{HTML}{DFE3DA}
\definecolor{pine}{HTML}{2C5C4C}    % structural accent
\definecolor{pinesoft}{HTML}{DCE6DF}
\definecolor{rust}{HTML}{A2432B}    % semantic: worse / caution
\definecolor{slate}{HTML}{40628C}   % semantic: better / down
\definecolor{amber}{HTML}{8A6A1F}   % findings

\color{ink}
% The web version sits on a tinted ground. On paper that costs ink and hurts
% print legibility, so the page stays white and the tint is reserved for boxes.
% Uncomment for an on-screen version that matches the artifact exactly:
% \pagecolor{paper}

%------------------------------------------------------------------------------
% TYPE
% The web report uses Spectral / Public Sans / IBM Plex Mono. Those are Google
% fonts; the Adobe Source trio below is the closest pairing that ships with
% Overleaf and compiles under pdfLaTeX with no uploads. Source Serif stands in
% for Spectral, Source Sans for Public Sans, Source Code Pro for Plex Mono.
%
% FOR THE EXACT WEB FONTS: switch the compiler to XeLaTeX or LuaLaTeX, upload
% the .ttf files, and replace this block with:
%   \usepackage{fontspec}
%   \setmainfont{Spectral}[Numbers=OldStyle]
%   \setsansfont{Public Sans}
%   \setmonofont{IBM Plex Mono}[Scale=0.88]
%   \renewcommand{\familydefault}{\sfdefault}
%------------------------------------------------------------------------------
\usepackage[T1]{fontenc}
\usepackage{sourceserifpro}                 % \rmfamily — headings
\usepackage[scale=0.96]{sourcesanspro}      % \sffamily — body
\usepackage[scale=0.86]{sourcecodepro}      % \ttfamily — data
\renewcommand{\familydefault}{\sfdefault}   % body is sans, as on the web
\usepackage{microtype}
% letterspace (microtype bundle) handles macros inside the argument;
% soul's \so does not and fails with 'Reconstruction failed'.
\usepackage{letterspace}

%------------------------------------------------------------------------------
% PAGE — single measured column, like the web report
%------------------------------------------------------------------------------
\usepackage[a4paper,margin=30mm,top=26mm,bottom=28mm]{geometry}
\linespread{1.10}
\setlength{\parskip}{0.62em}
\setlength{\parindent}{0pt}                  % web-style blocks, no indents
\usepackage{ragged2e}
\RaggedRight                                 % unjustified, as on screen

%------------------------------------------------------------------------------
% HEADINGS
%------------------------------------------------------------------------------
\usepackage{titlesec}
\titleformat{\section}
  {\rmfamily\bfseries\LARGE\color{ink}}{}{0pt}{}
\titlespacing*{\section}{0pt}{2.1em}{0.55em}
\titleformat{\subsection}
  {\rmfamily\bfseries\large\color{ink}}{\thesubsection\quad}{0pt}{}
\titlespacing*{\subsection}{0pt}{1.5em}{0.4em}
\titleformat{\paragraph}[runin]
  {\sffamily\bfseries\color{ink}}{}{0pt}{}[.]
\titlespacing*{\paragraph}{0pt}{0.9em}{0.5em}

% "eyebrow" — the small uppercase mono label above a section on the web
\newcommand{\eyebrow}[1]{%
  \par\vspace{0.4em}%
  {\ttfamily\footnotesize\color{pine}\textls[120]{\MakeUppercase{#1}}}\par\vspace{-0.5em}}

% section with an eyebrow above it
\newcommand{\esection}[2]{\eyebrow{#1}\section{#2}}

%------------------------------------------------------------------------------
% BOXES
%------------------------------------------------------------------------------
\usepackage[most]{tcolorbox}
\tcbset{
  boxrule=0pt, arc=0.6mm, left=5mm, right=5mm, top=3.2mm, bottom=3.2mm,
  boxsep=0mm, colback=raised, enhanced, breakable
}
% research-question block — pine rule on the left
\newtcolorbox{rqbox}{borderline west={2.4pt}{0pt}{pine}}
% finding — amber rule across the top, the headline result of a section
\newtcolorbox{findingbox}[1]{%
  borderline north={2.4pt}{0pt}{amber},
  title={\ttfamily\footnotesize\color{amber}\textls[120]{\MakeUppercase{#1}}},
  coltitle=amber, fonttitle=\ttfamily, attach title to upper=\par\vspace{1.5mm},
  colbacktitle=raised, titlerule=0pt, boxed title style={empty}}
% caution — rust rule on the left, for competing explanations
\newtcolorbox{cautionbox}[1]{%
  borderline west={2.4pt}{0pt}{rust},
  title={\sffamily\bfseries\color{rust}#1},
  coltitle=rust, attach title to upper=\par\vspace{1.5mm},
  colbacktitle=raised, titlerule=0pt, boxed title style={empty}}
% constructive — pine rule on the left
\newtcolorbox{goodbox}[1]{%
  borderline west={2.4pt}{0pt}{pine},
  title={\sffamily\bfseries\color{pine}#1},
  coltitle=pine, attach title to upper=\par\vspace{1.5mm},
  colbacktitle=raised, titlerule=0pt, boxed title style={empty}}

%------------------------------------------------------------------------------
% TABLES — booktabs rules, mono figures, tinted rows for significant results
%------------------------------------------------------------------------------
\usepackage{booktabs}
\usepackage{tabularx}
\usepackage{siunitx}
\sisetup{group-separator={,},detect-all=true}
\renewcommand{\arraystretch}{1.22}
\setlength{\aboverulesep}{0pt}\setlength{\belowrulesep}{0pt}
\arrayrulecolor{rule}

% column type for numeric cells: mono, right-aligned, tabular figures
\newcolumntype{N}{>{\ttfamily}r}
% tint a row whose result reached significance
\newcommand{\sigrow}{\rowcolor{pinesoft}}
% inline emphasis inside tables
\newcommand{\sig}[1]{\textbf{#1}}
\newcommand{\up}[1]{\textcolor{rust}{#1}}
\newcommand{\dn}[1]{\textcolor{slate}{#1}}
\newcommand{\nil}[1]{\textcolor{muted}{#1}}

% table captions above, small, muted
\usepackage[font={small,color=muted},labelfont={bf,color=ink},
            justification=raggedright,singlelinecheck=false,skip=5pt]{caption}

%------------------------------------------------------------------------------
% FIGURES
%------------------------------------------------------------------------------
% `float` gives [H]: tables land exactly where written instead of drifting
% above the heading they belong to, which matches the web report.
\usepackage{float}
\usepackage{graphicx}
\graphicspath{{figures/}}

%------------------------------------------------------------------------------
% LISTS — web-like spacing
%------------------------------------------------------------------------------
\usepackage{enumitem}
\setlist[itemize]{leftmargin=1.2em,itemsep=0.45em,topsep=0.4em}
\setlist[enumerate]{leftmargin=1.5em,itemsep=0.55em,topsep=0.4em}
\setlist[description]{leftmargin=0pt,itemsep=0.5em,topsep=0.4em,
                      font=\ttfamily\footnotesize\color{pine}}

%------------------------------------------------------------------------------
% MISC
%------------------------------------------------------------------------------
\usepackage{amsmath}
\usepackage[hidelinks]{hyperref}
\hypersetup{colorlinks=true,linkcolor=pine,citecolor=pine,urlcolor=pine}
\newcommand{\evt}{\tau}
% \colorbox is unbreakable and overflows the measure on long identifiers;
% \allowbreak lets a line break after an inline code span instead.
\newcommand{\code}[1]{%
  {\ttfamily\small\colorbox{pinesoft}{\hspace{0.15em}#1\hspace{0.15em}}}\allowbreak}
% \codep — for long paths and identifiers. A \colorbox is a single unbreakable
% object, so a long one runs into the margin; this keeps the mono face and
% drops the tint so the line can break normally.
\newcommand{\codep}[1]{{\ttfamily\small #1}}

% running footer, mono and muted, like the web footer
\usepackage{fancyhdr}
\pagestyle{fancy}\fancyhf{}
\renewcommand{\headrulewidth}{0pt}
\fancyfoot[C]{\ttfamily\scriptsize\color{muted}\thepage}

% title block
\newcommand{\reportheader}[4]{%
  {\ttfamily\footnotesize\color{pine}\textls[120]{\MakeUppercase{#1}}}\par\vspace{1.1em}
  {\rmfamily\bfseries\fontsize{26}{29}\selectfont\color{ink}#2}\par\vspace{0.7em}
  {\rmfamily\large\color{ink}#3}\par\vspace{0.9em}
  {\ttfamily\scriptsize\color{muted}#4}\par\vspace{0.5em}
  {\color{rule}\rule{\linewidth}{0.5pt}}\par\vspace{1.2em}}
